% =====================================================================
\chapter{Conclusiones}\label{cap:conclusiones}
% Plantilla ITBA: máximo 2 carillas; indicar si se cumplieron los
% objetivos y si la hipótesis fue comprobada o refutada.
% Borrador del 21 oct. 2026 (pendiente de revisión de la Dra. Kruse): una
% conclusión por asunto, estado de cada objetivo y de la hipótesis. No agrega
% resultados nuevos: resume los de la sección 5 con el enfoque de la Discusión.
% =====================================================================

En este trabajo se evaluaron los efectos a largo plazo del consumo de una bebida endulzada con un extracto de stevia durante la etapa juvenil o la adultez, y se los comparó con los de la sacarosa, utilizada como control positivo por sus efectos ya descriptos. De los resultados obtenidos se desprenden las siguientes conclusiones.

\textbf{Consumo.} Los animales juveniles consumieron, en relación con su peso corporal, más del doble de solución endulzada que los adultos, sin una mayor preferencia por el sabor dulce, y la ingesta de ambos \resaltar{endulzantes} superó los valores diarios de referencia.

\textbf{Conducta tipo ansiosa.} El consumo de stevia, en cualquiera de las dos edades, se asoció a largo plazo con una mayor exploración de las zonas centrales del \textit{open field}, sin cambios en la actividad locomotora, es decir, con un perfil de tipo ansiolítico. La sacarosa, en cambio, aumentó la conducta tipo ansiosa solo cuando se consumió durante la etapa juvenil.

\textbf{Memoria de reconocimiento.} La stevia alteró la memoria de reconocimiento de objetos de manera dependiente de la edad: a largo plazo en los animales expuestos durante la etapa juvenil y a corto plazo en los expuestos durante la adultez, sin afectar la memoria consolidada. La sacarosa consumida durante la etapa juvenil alteró los tres tipos de memoria evaluados.

\textbf{Separación de patrones espaciales y neurogénesis.} La stevia no alteró la separación de patrones espaciales en ninguna de las dos edades. La sacarosa la deterioró en ambas edades en la configuración de mayor demanda, y también en la de menor demanda cuando se consumió durante la etapa juvenil, y redujo en ambas edades la proliferación celular y el número de neuroblastos en proliferación en el \resaltarb{GD}: dos expresiones de un mismo efecto sobre la neurogénesis hipocampal.

\textbf{Análisis multivariado.} Las métricas de la prueba SLR se organizaron en dos dimensiones casi independientes, la exploración de los objetos y la discriminación de la posición novedosa. Los animales expuestos a stevia se ubicaron junto a los controles en la dimensión de la discriminación, y el clasificador distinguió a cada \resaltar{endulzante} del control por encima del azar, aunque ningún grupo constituyó un perfil conductual separado del resto.

\section*{Cumplimiento de los objetivos}

Se cumplieron los objetivos de mínima referidos al registro del consumo, a la conducta tipo ansiosa, a la memoria de reconocimiento y a la separación de patrones espaciales. \resaltar{La} cuantificación de PCNA y DCX se completó en los grupos control y sacarosa, pero no en los grupos expuestos a stevia. Los objetivos de máxima se cumplieron sobre la prueba SLR: el PCA y el clustering de Ward mostraron que los animales no se agrupan espontáneamente según su exposición, y el Random Forest predijo el grupo de cada animal por encima del azar, con el índice de discriminación como la variable de mayor importancia.

\section*{Hipótesis}

La hipótesis se cumplió solo en parte. La stevia no reprodujo los efectos adversos de la sacarosa sobre la memoria consolidada y la separación de patrones espaciales, pero no resultó inocua: produjo, en ambas edades, un perfil de tipo ansiolítico y alteraciones de la memoria de reconocimiento, de menor magnitud que las de la sacarosa. Su efecto sobre la neurogénesis hipocampal queda por determinar.

En conclusión, el consumo de stevia en etapas tempranas de la vida o en la adultez no está exento de efectos a largo plazo sobre el comportamiento, aunque estos son más sutiles que los de la sacarosa. Estos resultados indican que la seguridad de la stevia como alternativa al azúcar, especialmente durante el desarrollo, requiere una evaluación más profunda.
